% ============================================================
%  系统开发工具基础 · 实验报告 LaTeX 模板
%  作者：曾奕豪
%  编译方式：XeLaTeX（需 ctex 宏包支持中文）
%  用法：每周复制本文件为 weekN.tex，仅修改周次、日期、内容即可
%  说明：本模板不依赖任何图片即可编译；
%        图片文件存在时自动插入，不存在时显示占位提示框。
% ============================================================

\documentclass[12pt,a4paper]{article}

% ---------- 中文字体与页面 ----------
% 中文支持：使用 ctex 宏包（编译用 XeLaTeX）。
% 若在 Linux/TeX Live 下缺少 ctex，可安装 texlive-lang-chinese。
\usepackage[UTF8]{ctex}
\usepackage[margin=2.5cm]{geometry}
\usepackage{graphicx}
\graphicspath{{../figures/}}   % 图片统一放在 paper/figures/ 下
\usepackage{float}
\usepackage{booktabs}
\usepackage{xcolor}
\usepackage{listings}
\usepackage{enumitem}
\usepackage[colorlinks=true,linkcolor=black,urlcolor=blue]{hyperref}

% ---------- 页眉 / 页脚 ----------
\usepackage{fancyhdr}
\pagestyle{fancy}
\fancyhf{}
\fancyhead[R]{实验报告 · 第 \Week 周}
\fancyfoot[C]{\thepage}
\renewcommand{\headrulewidth}{0.4pt}

% ---------- 代码高亮环境 ----------
\lstset{
  basicstyle=\ttfamily\small,
  keywordstyle=\color{blue!70!black}\bfseries,
  commentstyle=\color{gray!70!black}\itshape,
  stringstyle=\color{red!60!black},
  numbers=left,
  numberstyle=\tiny\color{gray},
  frame=single,
  breaklines=true,
  tabsize=4,
  captionpos=b,
  showstringspaces=false,
}

% ---------- 图片自动插入宏 ----------
% 题内图片：图注自动按「图一、图二...」编号。
%   \resetpic              每题开头调用一次，重置编号
%   \resultimg[宽度]{图片文件名}   图片存在→插入；不存在→显示占位框（不报错）
% 自定义图注（如 commit 截图）：
%   \resultimgcap[宽度]{图片文件名}{图注}
\newcounter{picnum}
\newcommand{\resetpic}{\setcounter{picnum}{0}}
\newcommand{\resultimg}[2][0.8\textwidth]{%
  \stepcounter{picnum}
  \IfFileExists{\detokenize{#2}}{%
    \begin{figure}[H]
      \centering
      \includegraphics[width=#1]{\detokenize{#2}}
      \caption{图\chinese{picnum}}
    \end{figure}%
  }{%
    \begin{center}
      \fbox{\parbox{0.8\textwidth}{\centering\itshape [图片待补充] \detokenize{#2}\\\small 图\chinese{picnum}}}
    \end{center}%
  }%
}
\newcommand{\resultimgcap}[3][0.8\textwidth]{%
  \IfFileExists{\detokenize{#2}}{%
    \begin{figure}[H]
      \centering
      \includegraphics[width=#1]{\detokenize{#2}}
      \caption{#3}
    \end{figure}%
  }{%
    \begin{center}
      \fbox{\parbox{0.8\textwidth}{\centering\itshape [图片待补充] \detokenize{#2}\\\small #3}}
    \end{center}%
  }%
}

% ---------- 文档信息（每周只需改这三处） ----------
\newcommand{\Week}{1}                        % 周次
\newcommand{\ReportTitle}{系统开发工具基础实验报告}
\newcommand{\ReportDate}{\today} % 日期，编译时自动取当天
\newcommand{\RepoLink}{https://github.com/hzydy00/Base-of-system-development-tools} % 仓库链接

% ============================================================
\begin{document}

% ---------- 封面 / 标题区 ----------
\begin{center}
  {\LARGE\bfseries \ReportTitle}\\[6pt]
  {\Large 第 \Week\ 周}\\[12pt]
  {\large 姓名：曾奕豪}\\
  {\large 课程：系统开发工具基础}\\[6pt]
  {\large \ReportDate}\\[16pt]
  {\large 仓库：\url{\RepoLink}}
\end{center}
\vspace{0.8em}
\hrule
\vspace{1.2em}

% ---------- 1. 课上内容与结果 ----------
\section{课上内容与结果}
% 开头先简述本周实验的主题、目标与涉及的工具 / 技术，再逐题列出课上内容。
% 课上完成 4 个不同主题的题目，每题按"题目要求 → 代码 → 运行结果"组织。
此处填写本周实验的主题、目标与涉及的工具 / 技术。

\subsection*{题目一}
\resetpic
% TODO: 简述题目要求。
\begin{lstlisting}
# TODO: 题目一代码汇总

\end{lstlisting}

\resultimg{week1_check1.png}
\resultimg{week1_check1_2.png}

\subsection*{题目二}
\resetpic
% TODO: 题目要求、代码、结果
\begin{lstlisting}
# TODO: 题目二代码汇总
\end{lstlisting}

\resultimg{week1_check2.png}
\resultimg{week1_check2_2.png}

\subsection*{题目三}
\resetpic
% TODO: 题目要求、代码、结果
\begin{lstlisting}
# TODO: 题目三代码汇总
\end{lstlisting}

\resultimg{week1_check3.png}
\resultimg{week1_check3_2.png}

\subsection*{题目四}
\resetpic
% TODO: 题目要求、代码、结果
\begin{lstlisting}
# TODO: 题目四代码汇总
\end{lstlisting}

\resultimg{week1_check4.png}
\resultimg{week1_check4_2.png}

\subsection*{课上阶段提交记录}
% 课上 4 题的 commit 记录截图（git log 或仓库 commits 列表）
\resultimgcap[0.9\textwidth]{week1_commit_kecheng.png}{课上阶段 commit 记录}

% ---------- 2. 课后练习与结果 ----------
\section{课后练习与结果}
% 说明：课后完成 >10 个实例，每个实例按"题目 → 代码 → 运行结果"组织。

\subsection*{实例 1}
\resetpic
\begin{lstlisting}
# TODO: 实例1代码汇总
\end{lstlisting}

\resultimg{week1_ex1.png}
\resultimg{week1_ex1_2.png}

\subsection*{实例 2}
\resetpic
% TODO: 继续补充实例 2 ~ 实例 10+，结构与实例 1 相同
\begin{lstlisting}
# TODO: 实例2代码汇总
\end{lstlisting}

\resultimg{week1_ex2.png}
\resultimg{week1_ex2_2.png}

\subsection*{课后阶段提交记录}
% 课后练习的 commit 记录截图
\resultimgcap[0.9\textwidth]{week1_commit_kehou.png}{课后阶段 commit 记录}

% ---------- 3. 解题感悟 ----------
\section{解题感悟}
% TODO: 记录做题过程中的思考、遇到的问题与解决方式、收获与总结。
此处填写解题感悟。



\end{document}
